\begin{pastexam}[label={px:ncku-114-prog-10}]{成大資訊聯招 114 學年度 第 10 題}
（複選；至少有二個正確答案）

{[4\%]} Considering the following red-black tree, where the node colored as white indicates a red node and the node colored as black indicates a black node.

\begin{center}%
\begin{tikzpicture}[ every node/.style={font=\sffamily}, v/.style={ellipse,draw,minimum width=9mm,minimum height=6.5mm,inner sep=0pt}]
\node[v,fill=black,text=white] (12) at (4.98,0.00) {\textbf{12}};
\node[v,fill=black,text=white] (7) at (2.26,-0.87) {\textbf{7}};
\node[v,fill=black,text=white] (42) at (7.39,-0.85) {\textbf{42}};
\node[v,fill=black,text=white] (3) at (0.89,-2.10) {\textbf{3}};
\node[v,fill=black,text=white] (9) at (3.73,-2.03) {\textbf{9}};
\node[v,fill=white] (29) at (6.02,-2.08) {29};
\node[v,fill=black,text=white] (57) at (8.87,-2.00) {\textbf{57}};
\node[v,fill=white] (2) at (0.00,-3.29) {2};
\node[v,fill=white] (5) at (1.53,-3.32) {5};
\node[v,fill=white] (8) at (2.90,-3.27) {8};
\node[v,fill=black,text=white] (18) at (4.97,-3.35) {\textbf{18}};
\node[v,fill=black,text=white] (38) at (6.68,-3.27) {\textbf{38}};
\node[v,fill=white] (45) at (8.06,-3.19) {45};
\node[v,fill=white] (22) at (5.81,-4.40) {22};
\draw[] (12) -- (7);
\draw[] (12) -- (42);
\draw[] (7) -- (3);
\draw[] (7) -- (9);
\draw[] (3) -- (2);
\draw[] (3) -- (5);
\draw[] (9) -- (8);
\draw[] (42) -- (29);
\draw[] (42) -- (57);
\draw[] (29) -- (18);
\draw[] (29) -- (38);
\draw[] (18) -- (22);
\draw[] (57) -- (45);
\end{tikzpicture}%
\end{center}
\smallskip

Which of the following statements are correct?
\begin{enumerate}[label=(\Alph*)]
\item After inserting 1 into the tree, the number of black nodes becomes 9.
\item After successively inserting 100 and 58 into the tree, the number of red nodes becomes 7.
\item The deletion of 42 from the tree may reduce the height of the tree.
\item After successively deleting 45 and 57 from the tree, the rank of the root is 3.
\item The deletion of 12 can be done via moving 9 to the root.
\end{enumerate}
\end{pastexam}

\begin{pastsoln}
答案：(A)、(B)、(C)、(D)、(E) 都對。 (C) 要用「以中序前驅取代」的刪除法才成立，見下面的說明。

原樹（圖中白色是紅、黑色是黑）：黑節點 $12, 7, 42, 3, 9, 57, 18, 38$ 共 $8$ 個，紅節點 $29, 2, 5, 8, 22, 45$ 共 $6$ 個。 每條根到
外部節點的路徑都經過 $3$ 個黑色指標（例如 $12 \to 7 \to 3 \to 2 \to$ NIL：$7$、$3$、NIL），根的 rank（黑高度）是 $3$；樹高（邊數）
是 $4$（$12 \to 42 \to 29 \to 18 \to 22$）。 插入、刪除照資結書第~10.3.4~節、第~10.3.7~節（虛擬碼~10.6、10.8）。

\begin{itemize}
\item (A) 對。 插入 $1$：$12 \to 7 \to 3 \to 2$，成為 $2$ 的左子節點，紅色。 父 $2$ 紅、叔叔 $5$ 紅：$2$、$5$ 改黑，祖父 $3$ 改紅；
$3$ 的父節點 $7$ 是黑的，結束。 黑節點少了 $3$、多了 $2$、$5$，共 $8 - 1 + 2 = 9$ 個。
\item (B) 對。 插入 $100$：$12 \to 42 \to 57$，成為 $57$ 的右子節點，父節點黑，直接塗紅，紅節點變 $7$ 個。 插入 $58$：
$12 \to 42 \to 57 \to 100$，成為 $100$ 的左子節點；父 $100$ 紅、叔叔 $45$ 紅：$100$、$45$ 改黑，$57$ 改紅；$57$ 的父節點 $42$ 黑，
結束。 紅節點是 $29, 2, 5, 8, 22, 57, 58$，共 $7$ 個。
\item (C) 對（「may」）。 $42$ 有兩個子節點，要用中序後繼（$45$）或中序前驅（$38$）取代它，再刪掉那個節點。
\begin{itemize}
\item 用後繼 $45$：$45$ 是紅色葉子，直接刪除，不用修正。 $12 \to 45 \to 29 \to 18 \to 22$ 還在，樹高仍是 $4$。
\item 用前驅 $38$：$38$ 是黑色葉子，刪掉後 $29$ 的右邊少一個黑，要修正。 兄弟 $18$ 是黑的，遠姪子（$18$ 的左子）是 NIL、近姪子
$22$ 是紅的：先把 $22$ 改黑、$18$ 改紅，在 $18$ 左旋；再讓 $22$ 取 $29$ 的顏色（紅），$29$、$18$ 改黑，在 $29$ 右旋。 結果
$38$（原 $42$ 的位置）的左子樹是 $22$（紅），子節點 $18$、$29$（黑）。 最深的路徑變成 $3$ 條邊，樹高從 $4$ 降為 $3$。
\end{itemize}
所以刪除 $42$ 有可能讓樹變矮。 (E) 本身就用中序前驅來刪除，出題者接受這種做法；只用後繼的話，(C) 不成立。
\item (D) 對。 先刪紅色葉子 $45$，不用修正。 再刪黑色葉子 $57$：$42$ 的右邊少一個黑，兄弟 $29$ 是紅的：$29$ 改黑、$42$ 改紅，
在 $42$ 右旋，$29$ 接到 $12$ 的右邊，$42$ 成為 $29$ 的右子節點，$38$ 改接到 $42$ 的左邊。 新的兄弟 $38$ 是黑的、兩個子節點都是
NIL：$38$ 改紅，問題移到 $42$；$42$ 是紅的，改黑結束。 結果右子樹是 $29$（黑），子節點 $18$（黑，右子 $22$ 紅）、$42$（黑，左子
$38$ 紅）。 每條路徑仍是 $3$ 個黑色指標（例如 $12 \to 29 \to 18 \to$ NIL），根的 rank 是 $3$。
\item (E) 對。 $12$ 的中序前驅是左子樹最大的鍵：$7 \to 9$，$9$ 沒有右子節點，就是 $9$。 把 $9$ 搬到根（根維持黑色），再刪掉
$9$ 原本的節點：它是黑的、只有一個紅色子節點 $8$，讓 $8$ 接上來並改黑即可。 結果
$9(7(3(2, 5), 8), 42(\ldots))$，每條路徑仍是 $3$ 個黑色指標，是合法的紅黑樹。
\end{itemize}
\end{pastsoln}
